% Part: applied-modal-logics
% Chapter: epistemic-logic
% Section: relational-models

\documentclass[../../../include/open-logic-section]{subfiles}

\begin{document}

\olfileid{aml}{el}{rel}

\olsection{اړيکيز مدلونه}

د معرفتي منطقونو بنسټيز معنايي مفهوم د عادي مودالي منطقونو له مفهوم سره
يو شان دے۔ اړيکيز مدلونه لا هم د نړۍو له يوه سټ او داسې ټاکنې څخه جوړ
دي چې معلوموي کوم !!{propositional
variable}s په کومو نړۍو کښې «رښتيا» ګڼل کېږي۔ که يوازې له يوه عامل
سره کار لرو، د معمول په څېر د لاسرسي يوه اړيکه لرو۔ خو د څو-عامله
معرفتي منطق لپاره دغه يوه اړيکه د لاسرسي د اړيکو په يوه سټ بدلېږي؛
د عامل‌نښو په سټ~$G$ کښې د هر~$a$ لپاره يوه اړيکه۔

يو \emph{اړيکيز مدل} د نړۍو له يوه سټ، د هغو تر منځ له دوه‌ځاييزو
لاسرسي اړيکو—د هر عامل لپاره له يوې—او داسې ټاکنې څخه جوړ دے چې
معلوموي کوم !!{propositional variable}s په کومو نړۍو کښې رښتيا دي۔

\begin{defn}
  د څو-عامله معرفتي ژبې يو \emph{مدل} درې‌ګونې
  $\mModel{M} = \tuple{W, R, V}$ ده، چې په دې کښې
  \begin{enumerate}
  \item $W$ د «نړۍو» يو نا تش سټ دے،
  \item د هر $a \in G$ لپاره ${R}_a$ پر~$W$ د لاسرسي دوه‌ځاييزه
  اړيکه ده، او
  \item $V$ داسې تابع ده چې هر !!{propositional
    variable}~$p$ ته د ممکنو نړۍو يو سټ $V(p)$ ټاکي۔
  \end{enumerate}
  که $R_a ww'$ وي، نو وايو چې $w'$ ته \emph{د a له خوا
    له}~$w$ \emph{څخه لاسرسے شته}۔ که $w \in V(p)$ وي، نو وايو چې
  $p$ په~$w$ کښې \emph{رښتيا دے}۔
\end{defn}

چاره د نورمال مودالي منطق د چارې په شان ده، يوازې د لاسرسي اړيکې
پکښې ډېرې دي۔ د يوه ټاکلي عامل د لاسرسي اړيکه عموماً د هغه د
معلوماتي حالتونو د ښودلو په توګه تفسيروو۔ د بېلګې په توګه،
$R_a ww'$ ډېر ځله داسې لولو چې $w'$ په~$w$ کښې د~$a$ له
معلوماتو سره سازګاره ده۔ يا په بله وينا، په~$w$ کښې عامل د
نړۍ~$w$ او نړۍ~$w'$ تر منځ توپير نۀ شي پېژندلے۔

\end{document}
